\chapter{The Direct3D 9/11/12 Backends}
\label{ch:direct3d-backends}

All three of CNA's Windows-only Direct3D backends share one development loop:
cross-compiled from this project's own Debian machine using the same MinGW-w64 toolchain
Chapter~\ref{ch:building-cna}'s Windows cross-build section already introduced, then run and
verified under Wine on a real GPU --- \texttt{D3D9} and \texttt{D3D11} through DXVK,
\texttt{D3D12} through a different translation layer, vkd3d-proton, in its own separate Wine
prefix. None of the three has been verified on real Windows hardware yet; that remains open
for all three alike, tracked as its own \texttt{needs\_human} task in each backend's plan.

\section{D3D9: the one backend that targets pixel authenticity, not just parity}

\texttt{D3D9} is architecturally different in ambition from every other backend in this book.
Its own documentation states the goal directly: ``pixel-for-pixel indistinguishability from
the original XNA~4.0 runtime itself, not just `renders plausibly.''' It runs Microsoft's own
vendored, verbatim XNA~4.0 Stock Effects HLSL, compiled through the real
\texttt{d3dcompiler\_47.dll} rather than a reimplementation of it.

\subsection{The oracle: a real XNA~4.0 runtime as ground truth}

A dedicated tool, \texttt{tools/xna-oracle/}, stands up a genuine, running XNA~4.0 runtime
under Wine --- real GAC assemblies, compiled by the real in-prefix \texttt{csc.exe} --- and
renders the same declarative scene file through both that real runtime and CNA's \texttt{D3D9}
backend, diffing the two resulting images at zero tolerance: an exact match, not an
approximate one. Both sides of the comparison run through the same DXVK translation layer on
the same GPU, which is precisely what makes the comparison a test of CNA's own rendering
correctness rather than of two different graphics stacks. This corpus keeps growing across
this book's own successive passes over this chapter --- 36 scenes as of an earlier pass, 39
real scene files on disk as of this book's own most recent, direct recount --- with zero known
divergence at every count checked so far.

\subsection{What is real}

Full device lifecycle via plain \texttt{Direct3DCreate9} (not D3D9Ex, a deliberate choice),
with genuine \texttt{DeviceLost} / \texttt{DeviceResetting} / \texttt{DeviceReset} handling; all
five XNA stock effects, with 61 of 66 compiled shader variants byte-identical to Microsoft's
own shipped bytecode (the remaining five differ only by HLSL compiler version, and are
separately proven equivalent through the oracle rather than merely assumed so); a real,
oracle- and mutation-verified \texttt{SpriteEffect.fx}, including the half-pixel offset
Direct3D~9 is notorious for --- caught before it could become a false-positive ``closed''
claim, because the offset's absence is invisible on a $1\times1$ test texture and only shows
up at a larger size; and \cnaclass{GraphicsProfile.Reach} / \texttt{.HiDef} enforcement that is
\emph{real} here, consulting an actual \texttt{D3DCAPS9} structure --- the only CNA backend
where that distinction reflects a genuine capability query rather than a decorative property,
as Chapter~\ref{ch:graphicsdevice} already noted. Render targets sampled back as ordinary
textures, and the \texttt{PreferPerPixelLighting}/specular gaps from
Chapter~\ref{ch:stock-effects}'s \cnaclass{GpuDrawParams} discussion, were both real, open
problems earlier in this backend's history and are now fixed --- the specular fix in
particular proved four of five \texttt{PixelLighting} bytecode variants pixel-perfect via the
oracle; the fifth, an untextured vertex-lit variant, remains permanently unreachable, blocked
by the same missing position-only vertex layout that limits the untextured lit bucket on
every backend.

\subsection{What is not, and the one caveat that applies to every number above}

Every \cnaclass{SurfaceFormat} besides \texttt{Color} remains unsupported (a shared,
cross-backend limitation from Chapter~\ref{ch:textures-rendertargets}, not a D3D9-specific
one). \texttt{Texture3D} has no oracle-scene coverage, but it does have a real, capability-gated
sub-volume \texttt{SetData()}/\texttt{GetData()} byte round trip: \texttt{D3D9\_Smoke} creates a
\(4\times4\times4\) volume only when \texttt{D3DCAPS9::MaxVolumeExtent}
is nonzero, writes 32 distinct RGBA bytes into the off-origin \(2\times2\times2\) box at
\((1,1,1)\), and requires an exact readback of that same box. The test proves D3D9's
\texttt{LockBox()} row- and slice-pitch copy in both directions; it does not prove an oracle
rendering path or automatic mip generation. \cnaclass{OcclusionQuery} is not
built for this backend at all, the same gap D3D12 carries; and
\cnaclass{SpriteSortMode.Immediate} / \texttt{.Texture} were confirmed not viable as
deterministic oracle scenes at all --- \texttt{Immediate} is not something a raster diff can
observe, and \texttt{Texture} sorts by an implementation-defined \texttt{GetHashCode()},
which no two implementations are obligated to agree on. The caveat that applies to every
result in this section: \texttt{D3DCAPS9} itself is synthesized by DXVK on this dev loop, not
reported by genuine XNA-era driver hardware --- real-Windows-hardware verification remains
open, as stated above.

\subsection{D3D9 render targets prove native limits before promising a draw path}

The D3D9 smoke test proves real render-target storage and recovery rather than treating an
off-screen bind as a no-crash smoke test. An 8x8 target with Depth24Stencil8 is bound, cleared
through the public device, and copied with D3D9's required
\texttt{GetRenderTargetData()}--system-memory-surface--\texttt{LockRect()} route; its own surface
must contain the exact requested color. A cube-face target has the same proof for one selected
face, followed by an independent back-buffer clear/readback after unbind.

MSAA is capability-gated by \texttt{IDirect3D9::CheckDeviceMultiSampleType()}, not assumed from
the requested count. When 4x is available, a multisampled target is cleared and unbound; D3D9's
\texttt{StretchRect()} resolve must make its separate sampleable texture read exactly
\((200,210,220,255)\). The test proves the resolved result, not merely the multisample surface.
One naming trap matters here: this is an internal device-capability check in the creation path,
not \cnaclass{GraphicsDevice}\texttt{::SupportsCapability}. D3D9 never overrides that public
method, so \texttt{MultiSampleAntiAliasing} still reports true even when
\texttt{CheckDeviceMultiSampleType()} rejects every count above one. The same mismatch affects
\texttt{AnisotropicFiltering}: sampler creation uses \texttt{D3DTEXF\_ANISOTROPIC} and
\texttt{D3DSAMP\_MAXANISOTROPY}, but neither the public query nor the setter consults and clamps
to \texttt{D3DCAPS9::MaxAnisotropy}.

Multiple render targets are also bounded by the real \texttt{D3DCAPS9::NumSimultaneousRTs}: two
targets receive one exact-color clear, unbinding restores the back buffer, and requesting one
more target than the device reports throws rather than silently truncating. This proves D3D9's
multi-attachment binding and clear semantics. It does not separately prove a shader with several
fragment outputs, so the result should not be generalized into a G-buffer draw claim.

\subsection{A worked example: one real oracle scene, quoted in full}

The oracle's own scene format is a small, declarative text file, not a compiled asset or a
piece of C\texttt{++} --- the simplest real scene in the corpus, \texttt{colored3d.scene},
fits in eleven lines and is worth quoting directly rather than only described:

\begin{lstlisting}
width=256
height=256
profile=HiDef
clearcolor=100,149,237,255
vertexcolor=true
lighting=false
primitive=TriangleList
vertex=-0.6,-0.6,0,255,0,0,255
vertex=0.0,0.7,0,0,255,0,255
vertex=0.6,-0.6,0,0,0,255,255
\end{lstlisting}

This declares exactly the scene Chapter~\ref{ch:stock-effects}'s own \cnaclass{BasicEffect}
coverage would predict from the flags alone: \texttt{vertexcolor=true} with
\texttt{lighting=false} is \cnaclass{BasicEffect}'s \texttt{VertexColorEnabled} path with no
lighting applied, drawing a single red/green/blue-cornered triangle (World/View/Projection all
implicitly identity) over a \texttt{CornflowerBlue} clear --- the well-known default XNA clear
color, matching \texttt{Color.CornflowerBlue}'s own real RGBA value exactly
$(100, 149, 237, 255)$. Both the real XNA~4.0 oracle runtime and \texttt{D3D9} render this
identical declarative scene independently, and the zero-tolerance diff checks two things a
screenshot comparison alone would not separate cleanly: the flat-shaded corners prove basic
vertex-color plumbing works at all, while the \emph{center} of the triangle --- genuinely
Gouraud-interpolated between three different corner colors, not any one of them --- proves the
interpolation itself is correct, not merely that three isolated sample points happen to match.
A more elaborate scene in the same corpus, \texttt{fog\_gradient\_quad}
(Chapter~\ref{ch:easygl-backend}'s own negative-\texttt{FogEnd} finding), uses the identical
file format with a handful of additional \texttt{fog*} keys --- the same declarative scheme
scales from this eleven-line triangle up to every one of the corpus's now-39 scenes.

\subsection{Two real bugs the oracle methodology itself found}

Both are worth naming because they demonstrate what a pixel-exact reference actually buys
you over ordinary unit testing. First, a sort-mode Z-clipping bug: \texttt{zFarPlane=1}
mapped a nonzero \cnaclass{SpriteBatch} layer depth outside Direct3D~9's own valid clip-space
Z range, silently clipping away any sprite with a nonzero layer depth entirely --- fixed by
using \texttt{zFarPlane=-1} instead. Second, and more interesting because it is not
\texttt{D3D9}-specific at all: \cnaclass{GraphicsProfile} never actually reached the real
device, because \texttt{Game}'s own \cnaclass{GraphicsDevice} was eagerly default-constructed
(hardcoded to \texttt{Reach}) before \cnaclass{GraphicsDeviceManager} even existed to apply a
game's requested profile --- meaning a game setting \texttt{GraphicsProfile = HiDef;
ApplyChanges();} had no path to the live device at all. This is a shared, cross-backend bug,
only \emph{discovered} while building \texttt{D3D9}'s own tests, because \texttt{D3D9} is the
one backend where the profile actually matters enough to notice.

\subsection{A worked example: why \texttt{zFarPlane} sign alone hid every nonzero-depth sprite}

The Z-clipping bug above is worth working through numerically, since ``\texttt{zFarPlane=1}
instead of \texttt{-1}'' sounds like a cosmetic sign flip until the actual clip-space
consequence is spelled out. \cnaclass{SpriteBatch}'s orthographic projection maps a
\texttt{layerDepth} value in $[0, 1]$ (XNA's own documented range: 0 = front, 1 = back) onto
Direct3D~9's own native clip-space Z range, which is $[0, 1]$ --- \emph{not} the
$[-1, 1]$ range OpenGL-family APIs use. Constructing that projection with
\texttt{zFarPlane=1} produces a near/far mapping where any \texttt{layerDepth} strictly
greater than what the incorrect matrix treats as its own valid far bound gets clipped by the
GPU's own fixed-function Z-clip test --- silently, with no error, exactly like a sprite drawn
behind the camera would be. A sprite at \texttt{layerDepth=0} (the common case, since most
2D games never set it explicitly) was never affected, which is precisely why this bug could
ship unnoticed: it only manifests for a game that actually uses \cnaclass{SpriteSortMode}'s
depth-based ordering with a nonzero layer value, a real but less common \cnaclass{SpriteBatch}
usage pattern than the default. The fix, \texttt{zFarPlane=-1}, corrects the matrix so the
full documented $[0,1]$ \texttt{layerDepth} range maps inside Direct3D~9's own valid clip
volume instead of partly outside it.

\section{Three resource-lifetime groups, not two, and why the distinction matters}
\label{sec:d3d11-lifetime-groups}

Before getting to what \texttt{D3D11} actually renders, one design decision recorded in this
backend's own planning document (\texttt{plan\_dx.md}) is worth stating explicitly, because
getting it wrong is exactly the kind of mistake that makes a hand-rolled Direct3D backend
fragile under two of its most routine events --- a window resize and a lost device. A plain
resize and device-removed recovery touch genuinely different resource sets, and this backend
tracks that as three separate groups rather than the two-group split (``device stuff'' versus
``everything else'') an earlier draft of the same plan itself started from and explicitly
corrected:

\begin{description}
  \item[Device lifetime.] \texttt{ID3D11Device}, \texttt{ID3D11DeviceContext}, the
        \texttt{IDXGIFactory2} chain, and a cached \texttt{allowTearingSupported\_} capability
        flag --- created once at startup, torn down only by full device-removed recovery.
  \item[Swap-chain lifetime.] The \texttt{IDXGISwapChain1} object itself --- also created once
        at startup (and recreated only alongside the device, on recovery), but a plain resize
        deliberately \emph{reuses the same object} via
        \texttt{IDXGISwapChain1::ResizeBuffers(\ldots{})} rather than destroying and
        recalling \texttt{CreateSwapChainForHwnd}.
  \item[Window-size lifetime.] The back-buffer \texttt{ID3D11RenderTargetView}, the
        depth-stencil texture and its \texttt{ID3D11DepthStencilView}, and the viewport ---
        unbound and released \emph{before} \texttt{ResizeBuffers}, then recreated
        \emph{after} it, on every resize and on device-removed recovery alike.
\end{description}

The payoff of keeping these three separate: an ordinary window resize --- by far the most
frequent of the two events in practice --- touches only the third group plus one
\texttt{ResizeBuffers} call on the second, and specifically must \emph{never} re-run the first
group's factory/tearing-capability query or destroy and recreate the swap-chain object itself.
Device-removed recovery, the rarer and more disruptive event, is the only path that tears down
and rebuilds all three groups together. Collapsing swap-chain lifetime into either of the other
two groups --- treating it as part of ``device stuff'' (recreating it on every resize, wastefully)
or as part of ``window-size stuff'' (never giving it its own recovery-only teardown point) --- is
precisely the shape of bug this three-way split exists to rule out by construction, not merely
by convention.

A second, smaller decision from the same design-decisions list is worth a brief, concrete
mention for what it reveals about porting a hand-rolled COM-based backend under a non-MSVC
toolchain: every \texttt{ID3D11*} / \texttt{IDXGI*} / \texttt{ID3D12*} interface pointer anywhere in
this backend uses \texttt{Microsoft::WRL::ComPtr<T>}, not a bare \texttt{Release()} call or a
project-local wrapper --- but that was a genuinely verified choice, not an assumed one. Since
this project's actual primary dev loop is MinGW-w64 on Debian, not real MSVC, the plan records a
real spike (device creation, \texttt{.As()}, \texttt{GetParent}, and specifically the
\texttt{ReleaseAndGetAddressOf()} re-population pattern a swap-chain resize needs) that had to
compile, link, and run correctly under Wine on this project's own machine before the decision
was finalized, precisely because leaks on resize, double-releases, and forgotten releases on a
partially-failed initialization are the default failure mode for a hand-rolled COM wrapper that
was never actually exercised end to end.

\section{D3D11: the first native, dependency-free Windows path}

\texttt{D3D11} is CNA's first backend built directly on native Direct3D~11, verified in the
same MinGW-w64-plus-Wine-plus-DXVK loop as \texttt{D3D9}. It proves out a real
\texttt{ID3D11Device} / \texttt{ID3D11DeviceContext} pipeline: buffers, textures, and render
targets (including MSAA and multiple render targets), state objects, all ten stock HLSL
shader variants across the five stock effects, a real \cnaclass{SpriteBatch}, and --- unique
to this backend and \texttt{D3D12} --- a runtime-\texttt{D3DCompile()}-backed
\texttt{ShaderEffect} path for custom shaders, already introduced in
Chapter~\ref{ch:shader-fx-gap}. All of this is pixel-verified via real GPU readback, not
merely ``returned \texttt{S\_OK}.''

Device-lost/removed recovery here is honestly scoped as \emph{detection-only}: the
device-removed check is real and correctly wired, but has genuinely never fired in this dev
loop, since no real device removal has occurred under Wine+DXVK testing --- full recovery
remains unverified, not because the code is wrong, but because the triggering condition has
never actually happened. One real, structural bug is worth recounting for what it reveals
about static linking: an early build hit an \texttt{undefined reference to Effect::Apply()}
link failure, traced to a genuine circular dependency (the D3D11 \cnaclass{SpriteBatch}
backend calls back into \texttt{Effect::Apply()}, which lives in the main CNA library, which
itself depends on the backend). MinGW's single-pass archive resolution never revisited the
already-searched \texttt{libCNA.a} to satisfy it. The fix was an explicit repeated
\texttt{target\_link\_libraries(\ldots{} CNA)} for the D3D11/D3D12 targets, deliberately
exploiting CMake's own documented support for a static-library link cycle --- and the bug was
confirmed genuinely pre-existing (via \texttt{git stash}, it reproduced identically on the
base commit), not something the D3D11 work itself introduced.

\subsection{The public presentation-mode call reaches D3D11 and D3D12, but neither scales yet}

Chapter~\ref{ch:game-loop}'s \texttt{NOXNA} \texttt{Presentation\allowbreak{}Mode} is not lost before it
reaches these backends: the public manager/device forwarding path passes the selected enum to each
backend. But the final handlers are
currently intentional no-ops. D3D11 explicitly discards \texttt{mode}; D3D12 has the same
empty handler. They create and present their swap chain at the real SDL window's pixel size,
not a separately rendered logical surface which could be letterboxed, cropped, or stretched.

Their \texttt{SetVirtualResolution(width,height)} methods only retain the supplied integers.
On D3D11, \texttt{GetViewportSize()} still queries the window's physical pixel size; on D3D12,
the stored virtual dimensions are returned to the caller but do not create a scaling pass.
Consequently setting \texttt{Letterbox}, \texttt{Overscan}, \texttt{Stretch},
\texttt{NativeBackBuffer}, or \texttt{FixedHeightDynamicWidth} produces no different
presentation geometry on either native Direct3D backend today. This is a real backend limitation,
not an XNA-compatibility promise: the enum itself is a CNA extension, and there is no test that
could honestly pixel-prove a mode-specific result until a logical render target plus final
composite/scaling path exists.

\subsection{A proof ladder for real D3D11 3D draws, not a blanket ``3D works'' claim}

The D3D11 plan's own 3D work is unusually useful as a verification example because it did not
stop at compiling a shader or observing that a draw call returned successfully.  It built a
small ladder of increasingly discriminating GPU-readback checks.  The first rung used only the
stride-16 \cnaclass{VertexPositionColor} path: a known blue clear was read back, an NDC-space
opaque-red triangle was then drawn through both indexed and non-indexed calls, and the same
region had to become red.  That proves the input layout, shaders, constant buffer, primitive
topology, rasterization, and readback path together; a stale clear or an unbound shader cannot
pass it accidentally.

The next rung separated texture sampling from vertex-color modulation.  A stride-20 textured
draw sampled a known texel exactly, through both indexed forms, while the stride-24 variant
multiplied a known vertex color through a white texture and read back that exact color.  The
stride-32 lighting path then exercised the real per-light constant buffer: its unlit control is
byte-exact, while its lit control must differ from both the unlit result and the clear color.
That deliberately avoids pretending a hand-derived GPU floating-point result is more stable
than it is, yet still proves that the lighting branch actually executed.

Later tests applied the same discipline to features whose happy paths look deceptively similar.
Fog draws the same geometry with the flag off (exact vertex red) and on at the exact end distance
(exact fog green); an environment-map fixture constrains the reflection to the cube's \texttt{-Z}
face and reads its unique color; and a skinned fixture supplies a genuine identity bone rather
than relying on an all-zero default matrix.  These are deliberately different falsification
cases, not redundant screenshots.  By the end of that series the D3D11 smoke executable had
147 checks, each exercised through real Wine plus DXVK GPU execution.  The result supports the
specific paths named here; it is not a license to extrapolate that any arbitrary, untested
D3D11 state combination must therefore be correct.

\section{An MRT set has to finish every target, not only the first}

An ordinary single \cnaclass{RenderTarget2D} owns a simple end-of-use rule. When it is unbound,
its \texttt{ResolveAnd\allowbreak{}Generate\allowbreak{}MipsEXT()} helper asks D3D11 to resolve the
multisampled draw attachment into the separate sampleable texture. If the target
has a mip chain, it then calls \texttt{GenerateMips()} on the shader-resource view. Thus the
texture a later sprite samples is the resolved, current image rather than the still-multisampled
attachment or an old mip level.

That rule was once silently absent for an \(N>1\) MRT bind. The single-target pointer could not
represent several targets, so the old \texttt{SetRenderTargets()} path bound all their RTVs with
\texttt{OMSetRenderTargets()} but did not call any target's finalization when the set was
replaced. A game could clear or draw to two MSAA targets successfully, then sample stale
single-sample resolve textures without an exception. The repair keeps a non-owning array plus a
count for the active MRT set. Both \texttt{SetRenderTargets()} and
\texttt{SetRenderTarget2D()} first run \texttt{Flush\allowbreak{}Pending\allowbreak{}MRT\allowbreak{}ResolveEXT()}, which visits every
stored target, invokes its \texttt{ResolveAnd\allowbreak{}Generate\allowbreak{}MipsEXT()}, clears the array entry, and only
then binds the next target or restores the back buffer. This matters for all three transitions:
MRT to back buffer, MRT directly to one target, and MRT to a different MRT set.

\texttt{D3D11\_Smoke} makes the error observable rather than only reading that loop. It binds
two distinct \(8\times8\), 4x-MSAA targets, clears the shared MRT pass to
\((200,30,40,255)\), unbinds it, and reads both targets' \emph{sampleable} textures directly.
Both must contain that exact color. A repair that resolved only target zero, or that merely left
the MSAA draw attachment valid, fails the second readback. A second case binds another two-target
set, clears it to \((5,6,7,255)\), and switches straight to an unrelated single target rather
than unbinding to null; its first prior target must still read the new value. That separates the
real state-transition repair from a narrower ``only the null unbind works'' patch.

The shared helper also calls \texttt{GenerateMips()}, but the scope of its proof should remain
honest. The suite separately reads an \(8\times8\) single-target mipmapped render target's
\(4\times4\) and \(2\times2\) levels after a solid \((200,90,10,255)\) clear, proving that the
normal unbind path creates non-garbage downstream levels. The \emph{combination} of \(N>1\)
MRT and \texttt{mipMap=true} uses the same helper and is therefore wired, but is not independently
pixel-tested yet. That is a verification boundary, not a claim that the two mechanisms are
magically proven together.

\section{D3D12: built on D3D11's own foundation, and an honest split verification story}

\texttt{D3D12} deliberately reuses \texttt{D3D11}'s own HLSL/DXBC bytecode, format and state
mapping tables, and constant-buffer layouts --- \texttt{D3DCommon} from
Chapter~\ref{ch:backend-architecture} is exactly this shared foundation. It proves out a real
device, command queues, descriptor heaps, command lists, fences with genuine two-frame
back-pressure, a real per-resource barrier-transition tracker, pipeline state objects with
runtime-settable state, root signatures, real per-slot dynamic sampler state, buffers, 2D/cube/%
3D textures, real render targets (including multiple render targets with real mip-chain
generation and device-queried MSAA), real occlusion queries, all ten stock shader variants
(reusing \texttt{D3D11}'s own compiled bytecode directly, not re-deriving it), a real
\cnaclass{SpriteBatch} with the same runtime-compiled custom-shader support as
\texttt{D3D11}, and a real, functionally-proven device-removed recovery path.

The single most important structural fact about this backend's own verification is worth
stating plainly, because it qualifies every number above: \textbf{the routine CTest suite for
this backend is entirely off-screen}. Real swap-chain creation and \texttt{Present()} through
a live window are proven, but only through a separate, manual diagnostic tool run through a
Proton-managed launch --- never through the routine, automated CTest run, because a full
Proton bootstrap was judged too heavy and slow for a normal test run in this dev loop.

\subsection{Backbuffer readback is not part of the shared D3D11 foundation}

The repeated phrase ``GPU readback'' in this chapter names several different mechanisms.
D3D11 has a real \texttt{ReadBackbuffer()} override: it always copies
\texttt{backBufferTexture\_} into a staging texture, honors \texttt{RowPitch}, and backs the
public \texttt{GraphicsDevice::GetBackBuffer\allowbreak Data()} route with exact windowed
clear/draw
tests. It does not redirect to a currently bound render target.

D3D12 has no corresponding override. A public backbuffer call reaches
\cnaclass{IGraphicsBackend}'s default \texttt{runtime\_error} on a windowed device, while
\texttt{HeadlessEXT} has no swapchain to read in the first place. The D3D12 smoke suite's many
exact GPU pixel assertions call its working render-target-resource readback helpers directly;
the shared public SpriteBatch example below likewise reads the off-screen render target, not
\texttt{GetBackBufferData()}. Those results prove the named draw/resource path but cannot be
quoted as evidence for a public backbuffer bridge that does not exist. The complete backend
matrix and FNA contrast are in \S\ref{sec:backend-readback-matrix}.

\subsection{Two frames in flight is a real primitive, not a promise that every call site is non-blocking}

Unlike the driver-managed D3D11 model, D3D12 makes the command allocator lifetime explicit.
CNA therefore creates exactly two \texttt{ID3D12CommandAllocator} objects
(\texttt{kFramesInFlight = 2}), one reusable graphics command list, one shared
\texttt{ID3D12Fence}, and one remembered fence value per allocator slot.  The essential
operation, \texttt{SignalAndWaitForFrameEXT(frameIndex)}, deliberately has a slightly
counter-intuitive ordering: it signals a \emph{new}, monotonically increasing value for the
frame just submitted, then waits only for the \emph{previous} value recorded for that same
slot before allowing its allocator to be reused.  Thus a call for slot~0 may submit value
\texttt{v2} while waiting for older \texttt{v0}; it must not wait for \texttt{v2} itself,
because doing so would erase the overlap the two slots exist to permit.

The backend's real smoke test makes that contract observable rather than treating fence
arithmetic as self-evident.  It invokes the primitive for slots 0, 1, and 0 again, verifies
that \texttt{v0 < v1 < v2}, and, on the second use of slot~0, requires
\texttt{GetCompletedValue()} to have reached \texttt{v0}.  It intentionally does \emph{not}
assert that \texttt{v2} has completed on return: \texttt{ID3D12CommandQueue::Signal()} is
asynchronous, so that would be a race, not a correctness condition.  A separate explicit
event wait proves eventual completion.  The test goes one step further by placing 48 real
64-MiB GPU copies ahead of the old fence and comparing that wait with an already-complete
control wait; the loaded wait is measurably slower, establishing that the
\texttt{WaitForSingleObject()} branch genuinely supplies back-pressure rather than merely
passing an already-satisfied condition.

There is an important present-tense limitation in the distinction.  The reusable
two-slot primitive is real and directly tested, but the current ordinary clear/draw helper
uses \texttt{ExecuteCommandListAndWaitEXT()}, its intentionally simpler synchronous sibling:
it submits a closed list, signals a fresh fence value, and waits for \emph{that} value before
returning.  \texttt{Present()} likewise uses that synchronous helper for its explicit
back-buffer transition before calling \texttt{IDXGISwapChain3::Present()}.  Consequently a
reader should understand ``two frames in flight'' here as tested infrastructure and the
correct allocator-reuse rule, not as a performance claim that all current CNA D3D12 drawing
automatically overlaps two CPU frames.  This is an honest staging choice: the off-screen
test paths prioritize a deterministic readback boundary, while the primitive preserves the
non-stalling policy a batched production present loop needs when that path is used.

The descriptor lifetime model has the same deliberately explicit character.  Render-target,
depth-stencil, shader-resource, and sampler descriptors come from four separate heaps; their
current allocators are fixed-capacity, monotonically advancing bump allocators, not a general
free-list.  A destroyed resource does not return its descriptor slot during the device's
lifetime, and exhaustion throws a named \texttt{std::runtime\_error} instead of silently
aliasing another resource's view.  The capacities in the live backend --- 64 RTV, 8 DSV,
64 CBV/SRV/UAV, and 16 sampler descriptors --- are therefore engineering limits of this
implementation, not Direct3D~12 limits.  The RTV heap has already been raised several times
when genuine smoke-test render targets, including a six-face cube target, consumed its earlier
budget.  That history is a useful warning for code that creates and destroys a large number
of transient CNA render targets: it should not mistake the absence of a descriptor free-list
for a GPU-memory leak, but it also cannot assume unlimited descriptor churn in one device
instance.

\subsection{A genuine architecture mismatch, not a CNA bug}

Swap-chain creation under \emph{plain} Wine (as opposed to a Proton-managed launch) crashes
outright, reproduced twice with a full symbolized backtrace pointing to a null-pointer read
inside Wine's own \texttt{dxgi.dll}. The root cause, once found, is genuinely external to
CNA: a mismatch between Debian's own system \texttt{dxgi.dll} and vkd3d-proton's separately
overridden \texttt{d3d12.dll} --- two DLLs that were never meant to be paired this way. Once
launched through Proton properly, which supplies the matched DLL pair vkd3d-proton actually
expects, swap-chain creation and a genuine ten-frame clear-and-present loop through a real
window both succeed cleanly, verified twice with zero crashes.

\subsection{Explicitly corrected claims}

\texttt{D3D12}'s own documentation is unusually direct about its own earlier revisions being
stale: runtime-settable \cnaclass{BlendState} / \cnaclass{DepthStencilState}/%
\cnaclass{RasterizerState} feeding real pipeline-state objects, sixteen-slot dynamic sampler
state, real public render-target and render-target-cube backends (including multiple render
targets), real \cnaclass{Texture3D}, and real \cnaclass{TextureCube::GetData()} were all
previously described as open gaps in earlier documentation passes and are now all real and
closed --- with the documentation's own closing line telling a reader to re-verify any
\emph{other} specific claim against the project's own live tracking file before trusting it,
rather than this book's snapshot alone.

\section{One public \texttt{ShaderEffect} API, two fixed Direct3D contracts}

Chapter~\ref{ch:shader-fx-gap} now draws an important boundary: D3D11 and D3D12 really do
compile custom HLSL at runtime, but their implementations are \cnaclass{SpriteBatch} custom
shader facilities, not general replacements for XNA's compiled \texttt{Effect}. The two
backends accept the same constructor strings, compile the same \texttt{main} entry points as
\texttt{vs\_5\_0} and \texttt{ps\_5\_0}, and expose the same public setters. Under that shared
surface, both require one exact 32-byte vertex:

\begin{longtable}{@{}p{0.22\textwidth}p{0.18\textwidth}p{0.46\textwidth}@{}}
\toprule
\textbf{Semantic} & \textbf{Byte offset} & \textbf{Required HLSL type} \\
\midrule
\endhead
\texttt{POSITION0} & 0  & \texttt{float2}: pixel-space X and Y \\
\texttt{TEXCOORD0} & 8  & \texttt{float2}: texture U and V \\
\texttt{COLOR0}    & 16 & \texttt{float4}: sprite tint \\
\bottomrule
\end{longtable}

This is \cnaclass{SpriteBatch}'s own vertex, not an inferred convention. D3D11 creates a
three-element \texttt{ID3D11InputLayout} from the compiled vertex blob; D3D12 embeds the same
three elements in the custom effect's pipeline-state object. A shader declaring
\texttt{float3 POSITION}, a tangent, or a game-specific stride therefore cannot be made
compatible by supplying a matching \cnaclass{VertexDeclaration}: unlike EasyGL's separately
tested general-3D path, neither Direct3D backend reads
\texttt{GpuDrawParams::customEffectBackend} from its 3D draw dispatch.

\subsection{The names are documentation; the byte offsets are the API}

Both implementations store public uniform writes in the same 128-byte block. The caller's
\texttt{name} argument is deliberately ignored. What selects a value is which setter was
called:

\begin{longtable}{@{}p{0.19\textwidth}p{0.25\textwidth}p{0.48\textwidth}@{}}
\toprule
\textbf{Bytes} & \textbf{Writer} & \textbf{Meaning} \\
\midrule
\endhead
0--15 &
backend only &
\texttt{vpSize}; \cnaclass{SpriteBatch} writes width and height before binding the effect. \\
16--79 &
\texttt{SetUniformMat4} &
one 64-byte matrix, regardless of the supplied name. \\
80--95 &
\texttt{SetUniformVec2/3/4} &
one overlapping vector slot; a later vector setter overwrites the components it supplies. \\
96--99 &
\texttt{SetUniformFloat/Int} &
one scalar slot. The integer overload stores a floating-point conversion, not integer bits. \\
100--127 &
none &
reserved padding in the current implementation. \\
\bottomrule
\end{longtable}

There are two further consequences that the common public surface does not make obvious.
\texttt{SetUniformFloatArray()} and \texttt{SetUniformVec2Array()} inherit
\cnaclass{IEffectBackend}'s no-op defaults on both backends. So do all three
\texttt{SetTexture()} backend hooks. The sprite's primary resource still works because
\cnaclass{SpriteBatch} itself binds its texture and sampler to \texttt{t0}/\texttt{s0}; a
second 2D texture, cube map, or volume texture set through \cnaclass{ShaderEffect} does not.
Effect authors should therefore treat the layout above as a tiny fixed ABI, not as
name-reflected HLSL constants.

\subsection{Where D3D11 and D3D12 stop being the same}

D3D11 can bind the compiled vertex shader, pixel shader, input layout, and dynamic constant
buffer as separate context state. Every \texttt{Bind()} performs a write-discard map, copies
the 128 bytes, and attaches the buffer to \texttt{b0} for both shader stages. The native map
mode is \texttt{D3D11\_MAP\_WRITE\_DISCARD}.

D3D12 has no equivalent independent shader bind. Its effect constructor must build a complete
\texttt{ID3D12PipelineState} up front, using the same cached root-signature shape as the stock
sprite path: one CBV at \texttt{b0}, one SRV table at \texttt{t0}, and one sampler table at
\texttt{s0}. The 128-byte constant buffer lives in a persistently mapped upload resource;
\texttt{Bind()} only copies the current values, while
\texttt{D3D12SpriteBatchBackend::FlushBatch()} binds the resource and the prebuilt PSO inside
its command list.

That PSO also makes D3D12's boundary stricter and directly visible in source: it is compiled
for one \texttt{DXGI\_FORMAT\_R8G8B8A8\_UNORM}, single-sample color target, with blending,
depth, and stencil disabled and culling set to none. Those choices match the real custom-effect
pixel test. They do not form a dynamic promise that the same object can be reused for an MRT
set, an MSAA target, another color format, or caller-selected pipeline state; those dimensions
would require separate PSOs and a cache key that the current effect backend does not build.

\subsection{A worked example shared by the D3D11 and D3D12 pixel tests}

The two smoke tests use byte-for-byte equivalent HLSL to turn a solid red texture into exact
cyan. The vertex shader's \texttt{vpSize} field is at byte zero because the backend fills it;
the four padding vectors advance the next public vector slot to byte 80:

\begin{lstlisting}
const char* vertexHlsl = R"(
struct VSIn  { float2 pos : POSITION0;
               float2 uv  : TEXCOORD0;
               float4 col : COLOR0; };
struct VSOut { float4 pos : SV_Position;
               float2 uv  : TEXCOORD0;
               float4 col : TEXCOORD1; };
cbuffer CB : register(b0) {
    float4 vpSize;
    float4 pad1[4];
    float4 uColor;
    float4 uFloat0;
};
VSOut main(VSIn input) {
    VSOut output;
    float2 ndc = (input.pos / vpSize.xy) * 2.0 - 1.0;
    output.pos = float4(ndc.x, -ndc.y, 0.0, 1.0);
    output.uv  = input.uv;
    output.col = input.col;
    return output;
})";

const char* pixelHlsl = R"(
Texture2D texSampler : register(t0);
SamplerState texSamplerSampler : register(s0);
struct PSIn { float4 pos : SV_Position;
              float2 uv  : TEXCOORD0;
              float4 col : TEXCOORD1; };
float4 main(PSIn input) : SV_Target {
    float4 source = texSampler.Sample(texSamplerSampler, input.uv);
    return float4(1.0 - source.rgb, 1.0);
})";

ShaderEffect invert(device, vertexHlsl, pixelHlsl);
if (!invert.IsEffectValid()) {
    throw std::runtime_error("custom HLSL did not compile");
}

SamplerState pointClamp = SamplerState::PointClamp;
SpriteBatch batch(device);
batch.Begin(SpriteSortMode::Deferred, BlendState::Opaque,
            &pointClamp, nullptr, nullptr, &invert);
batch.Draw(redTexture, destination, Color::White);
batch.End();
\end{lstlisting}

D3D11 runs this through a windowed Wine/DXVK device and reads the sprite's top-left region.
D3D12 runs the same public \cnaclass{GraphicsDevice}/\cnaclass{SpriteBatch} sequence against a
\texttt{HeadlessEXT} off-screen device, then reads the render target. In both cases solid
\((255,0,0,255)\) must become exactly \((0,255,255,255)\). That result proves more than
successful compilation: the fixed input layout, automatic viewport slot, \texttt{b0}
constant-buffer/root binding, \texttt{t0}/\texttt{s0} resource path, custom pixel shader, and
GPU readback all participated in one observable public-API draw.

\section{A reopened feature, and a public routing gap above the finished mechanisms}

\texttt{D3DCommon}'s shared HLSL/DXBC and state-mapping tables, named above, make it easy to
assume the two backends are near-identical beneath that layer. A late pair of tasks in the
Direct3D~11/12 implementation plan --- opened well after both backends were otherwise
considered feature-complete, once a full feature-matrix percentage audit against \texttt{EasyGL}
turned up a handful of remaining gaps --- shows a case where that assumption holds
(\cnaclass{RenderTargetCube} MSAA) and one where it genuinely does not (mip-chain generation).
The mechanisms described below are real, but a later interface-default audit found that the
book's former conclusion was too broad: both test suites invoke the cube backend's
\texttt{BindAsRenderTargetFace()} and \texttt{UnbindAsRenderTarget()} directly.
\cnaclass{GraphicsDevice}'s public cube-target route invokes the first but, on D3D11/12,
currently never invokes the second. The following subsections therefore describe working
backend algorithms with a missing public finalization edge, not end-to-end public features.

\subsection{\cnaclass{RenderTargetCube} MSAA: a deliberate exclusion, reopened on request}

Multisampled render-target cubes were, at one point, a matched, deliberate scope exclusion on
both backends, not a silent oversight --- the project owner explicitly asked for it to be
reopened alongside the rest of a late feature-matrix sweep. The two implementations that
resulted are a real worked example of the same underlying Direct3D constraint solved twice,
independently, once per API generation. Neither \texttt{D3D11\_RESOURCE\_MISC\_TEXTURECUBE} nor
a \texttt{D3D11\_SRV\_DIMENSION\_TEXTURECUBE} view can ever be multisampled, and \texttt{D3D12}'s
equivalent \texttt{D3D12\_SRV\_DIMENSION\_TEXTURECUBE} has no multisampled variant either --- so
on both backends, the MSAA color resource becomes a plain, non-cube six-slice
\texttt{Texture2DMSArray} used only as a render-target view, while a second, single-sample,
genuinely cube-flagged resource is resolved into via \texttt{ResolveSubresource()} on unbind, and
it is that resolved resource, not the MSAA one, that the real shader-resource view always
targets. \texttt{D3D12}'s own version of that same resolve additionally requires two explicit
resource-barrier transitions (\texttt{RESOLVE\_SOURCE} on the color resource,
\texttt{RESOLVE\_DEST} on the resolve target) that \texttt{D3D11} does not need at all --- its
context-level \texttt{ResolveSubresource()} call handles the equivalent state transition
implicitly. Both backends resolve only the currently active face on unbind, the same
one-face-at-a-time convention their mip-chain handling (below) also follows. Real proof on both
backends: an $8\times8$ cube face requested at 4x MSAA, cleared, resolved on unbind, and read
back exactly matching, with an explicit check that \texttt{GetMultiSampleCount()} reports the
real requested sample count rather than a silent single-sample fallback.

\subsection{Mip-chain generation: one native GPU call versus an explicit CPU round trip}

Generating a \cnaclass{RenderTargetCube}'s mip chain after a face is drawn to is where the two
backends' shared foundation runs out. \texttt{D3D11} calls
\texttt{ID3D11DeviceContext::GenerateMips()} directly on the cube's own shader-resource view ---
a single native, hardware-accelerated call with no per-face argument at all, since the view it
operates on spans the whole six-face resource. Calling it after any one face is drawn
regenerates the entire cube's mip chain from each face's own current level-0 content in one GPU
call; this is harmless for the five faces that were not just drawn to (their level-0 content did
not change, so recomputing their mips from it is a no-op in effect) and correct for the one that
was. \texttt{D3D12} has no equivalent convenience API for generating mips on an arbitrary format
at all, so \texttt{D3D12RenderTargetCubeBackend::GenerateMipsEXT()} implements the whole chain by
hand: for each mip level, it reads the previous level's pixels back from the GPU for the one
active face only (\texttt{ReadbackSubresourceRGBA8()}), box-filter-downsamples them on the CPU
(\texttt{BoxFilterDownsample()}), and uploads the result into the next level
(\texttt{UploadSubresourceRGBA8()}) --- an explicit, per-level CPU round trip, deliberately scoped
to only the face that was actually just drawn, since repeating that CPU work for all six faces on
every unbind would be needless cost for five faces whose content had not changed. The two
mechanisms reach the same visible result --- a correctly downsampled mip chain for the face that
was drawn --- through genuinely different means: one hardware call operating on the whole
resource at once, and one explicit, face-scoped software loop working level by level. It is a
concrete illustration of what ``\texttt{D3DCommon} is the shared foundation'' does and does not
promise: shared bytecode and state-mapping tables, not identical implementation strategy for
every individual feature, and a reader porting code that depends on either backend's own timing
or cost characteristics for this specific operation should not assume the other backend behaves
the same way underneath.

\subsection{Why the direct smoke proofs do not establish the public feature}

The missing edge is visible only by following the public call chain. The common cube-target
method on \cnaclass{IGraphicsBackend} binds a non-null cube directly and delegates a null cube
to the 2D-null path; it never calls the outgoing cube's unbind callback. Both Direct3D
backends inherit that default and track only a current 2D target. D3D11's 2D-null branch
assumes that target restored the backbuffer as part of its unbind, so after a public cube bind
it does not even restore the backbuffer. D3D12 restores the backbuffer unconditionally, but
still skips the cube's resolve/mipmap callback. Switching directly between cube faces has the
same missing-finalization shape on both.

The smoke checks cited above deliberately bypass this chain. They create a backend cube, call
\texttt{BindAsRenderTargetFace()}, clear it, and then call
\texttt{UnbindAsRenderTarget()} themselves before reading the resolved/mip resource. Their
exact pixels prove that each algorithm works when called, but cannot catch an absent caller.
No D3D11 or D3D12 example was found that performs the equivalent sequence through
\cnaclass{GraphicsDevice}'s public cube-target overload followed by the public null-target
overload. Until backend-level cube tracking or explicit overrides close that edge, public
D3D11 cube rendering is unsafe even at single-sample level~0 because later backbuffer work may
still hit the cube; public D3D12 single-sample level~0 survives, but cube MSAA resolve and
generated mips do not.

\section[\texttt{HeadlessEXT}: letting three D3D12 tests join the routine suite]{\texttt{HeadlessEXT}: the NOXNA property that let three D3D12 tests join the
routine suite}

Three specific gaps --- D3D12's own \cnaclass{SpriteFont} rendering, its \cnaclass{Model} runtime
orchestration, and half of \texttt{Texture2D.FromStream} / \texttt{SaveAsPng} --- stayed open
noticeably longer than every comparable D3D11 row, and a later planning pass traced all three to
one shared architectural root cause rather than three unrelated gaps: each needs a real
\cnaclass{GraphicsDevice}, and \cnaclass{GraphicsDevice}'s own constructor unconditionally
initialized SDL video and created a real window for every backend except \texttt{HEADLESS}/%
\texttt{SOFTWARE}. For \texttt{D3D12} specifically, a real window means a real swap chain, which
on this project's own Linux dev loop only works through the heavy, Proton-managed launcher
Section~\ref{sec:d3d11-lifetime-groups}'s sibling section above already introduced --- not the
plain Wine the routine \texttt{ctest} run uses. Three tests genuinely needing a
\cnaclass{GraphicsDevice} at all, on a backend whose only route to one was a launcher too slow for
routine CI, is a real architectural dead end, not a missing-test-case gap ordinary engineering
effort could close directly.

The fix actually shipped, rather than merely proposed: \cnaclass{PresentationParameters} gained a
NOXNA \texttt{HeadlessEXT} property, read directly from the real header:

\begin{lstlisting}
// CNA extension. When true, GraphicsDevice creates no SDL window and never
// initialises SDL's video subsystem, so it can run with no display server at
// all -- the same thing the Headless and Software backends do
// unconditionally, but as an opt-in for a backend that normally does want a
// window. Only backends that can genuinely operate without a swap chain
// support this. Today that means D3D12 ... Present() is not meaningful
// without a swap chain; an off-screen device is for rendering into render
// targets and reading the result back (tests, server-side rendering,
// thumbnail generation), not for putting frames on screen.
NOXNA bool getHeadlessEXTProperty() const;
NOXNA void setHeadlessEXTProperty(bool value);
\end{lstlisting}

Defaulting to \texttt{false} keeps production behavior byte-identical for every existing backend
and caller; only \texttt{D3D12} honors it today, since its own backend already treated a null
window as a real off-screen mode internally, while \texttt{D3D11}'s constructor always creates a
swap chain unconditionally and \texttt{EasyGL}'s GL context is inseparable from a window --- both
throw if a caller sets \texttt{HeadlessEXT} against them. With the property in place, all three
previously-stuck rows landed in the \emph{routine}, plain-Wine \texttt{D3D12} \texttt{ctest} run
like every other check, no Proton launcher required. \texttt{HeadlessEXT} is deliberately framed
in its own header comment as more than a test convenience: an off-screen \cnaclass{GraphicsDevice}
is a real, general capability --- server-side rendering and thumbnail generation being the two
named examples --- that happens to have been motivated by closing a testing gap first.

\subsection{Four real bugs a headless \texttt{GraphicsDevice} immediately found on D3D12}

Landing that first real, routine \cnaclass{GraphicsDevice}-driven test coverage on \texttt{D3D12}
did what this book has found repeatedly elsewhere: it surfaced bugs no amount of reasoning about
the code in isolation had caught, because nothing had ever driven the shared
\cnaclass{GraphicsDevice} path against this specific backend before. \cnaclass{SpriteBatch}
silently ignored \texttt{SetSamplerFilter} / \texttt{SetSamplerAddressMode} entirely, drawing with
whatever sampler state an unrelated, prior 3D draw happened to leave bound rather than the one the
caller actually requested. Render targets never bound their own depth-stencil view at all, so a
render target constructed \emph{with} a depth buffer silently behaved as though it had none ---
every depth test against it was inert, and a depth/stencil \texttt{Clear} had nothing real to
write to. All six combo \texttt{Clear*} variants (color-plus-depth, color-plus-stencil, and so on)
were still unimplemented throws. Most strikingly, \texttt{SetDepthTestEnabled} /
\texttt{SetDepthWriteEnabled} / \texttt{SetBlendEnabled} were also still unimplemented throws on
this backend --- meaning any real game calling any of these three ordinary, load-bearing
\cnaclass{GraphicsDevice} methods against \texttt{D3D12} did not render incorrectly, it
\emph{crashed} outright, found only because the \cnaclass{Model} test was the first thing in the
project ever to drive that specific call path against this backend. All four are now fixed; the
shape of the finding is the more durable lesson than any one bug: a genuinely new, real caller of
an already-``complete''-looking code path is one of this project's most reliable bug-finding
techniques, confirmed again here.
